\documentclass{article}

\usepackage{fancyhdr}
\usepackage{amsmath}
\allowdisplaybreaks
\usepackage{amsthm}
\usepackage{amsfonts}
\usepackage{tikz}
\usepackage[plain]{algorithm}
\usepackage{algpseudocode}
\usepackage{enumitem}
\usepackage{tikz}
\usepackage{pgfplots}
\pgfplotsset{compat=1.18}
\usepgfplotslibrary{groupplots}
\usepackage{comment}

\usetikzlibrary{automata,positioning}

%
% Basic Document Settings
%

\topmargin=-0.45in
\evensidemargin=0in
\oddsidemargin=0in
\textwidth=6.5in
\textheight=9.0in
\headsep=0.25in

\linespread{1.1}

\pagestyle{fancy}
\lhead{\hmwkAuthorName}
\chead{} % empty
\rhead{\hmwkHeaderTitle}
\cfoot{\thepage}

\renewcommand\headrulewidth{0.4pt}
\renewcommand\footrulewidth{0pt}

\setlength\parindent{0pt}

%
% Homework Details
%

\newcommand{\hmwkTitle}{Homework 3}
\newcommand{\hmwkDueDate}{Monday, January 26, 2026}
\newcommand{\hmwkClass}{ICMA 216 Calculus IIIA}
\newcommand{\hmwkClassTime}{Section 1}
\newcommand{\hmwkClassInstructor}{Ruth J. Skulkhu}
\newcommand{\hmwkAuthorName}{\textbf{Jiraroj Wiruchpongsanon (6781617)}}
\newcommand{\hmwkHeaderTitle}{ICMA 216 Calculus IIIA: Homework 3}

%
% Title Page
%

\title{
    \vspace{2in}
    \textmd{\textbf{\hmwkClass:\ \hmwkTitle}}\\
    \normalsize\vspace{0.1in}\small{Due\ on\ \hmwkDueDate}\\
    \vspace{0.1in}\large{\textit{\hmwkClassInstructor\ \hmwkClassTime}}
    \vspace{3in}
}

\author{\hmwkAuthorName}
\date{}

\renewcommand{\part}[1]{\textbf{\large Part \Alph{partCounter}}\stepcounter{partCounter}\\}

%
% Helper Commands
%

\newcounter{partCounter}
\newcommand{\solution}{\textbf{\large Solution}}

\newtheorem*{theorem}{Theorem}

\theoremstyle{remark}
\newtheorem*{remark}{Remark}

\begin{document}

\maketitle
\newpage

\section*{Problem 1 (11.3 Dot product)}
Use dot product to find the angle that vector $v = -\sqrt{3} i + j$ makes with the vector $w = i$, (positive $x$-axis).

\medskip
\solution

\vspace{6em}

\hrulefill

\section*{Problem 2 (11.3 Dot product)}
\begin{enumerate}[label=(\alph*)]
    \item Show that $(2,1,6)$, $(4,7,9)$ and $(11,7,-12)$ are the vertices of a right triangle.
    \item Find the area of the triangle.
\end{enumerate}

\medskip
\solution

\vspace{10em}

\hrulefill

\section*{Problem 3 (11.3 Dot product)}
Given the points $P_1 = (2, 1, 6)$, $P_2 = (4, 7, 9)$, and $P_3 = (11, 7, -12)$.
\begin{enumerate}[label=(\alph*)]
    \item Find the vector with length $2$ and has opposite direction to the vector $\overrightarrow{P_1P_2}$.
    \item Show that the vector $\overrightarrow{P_1P_2}$ is orthogonal (perpendicular) to the vector $\overrightarrow{P_1P_3}$.
\end{enumerate}

\medskip
\solution

\vspace{10em}

\hrulefill

\section*{Problem 4 (11.3 Dot product)}
True-False Determine whether the statement is true or false. Explain your answer.
\begin{enumerate}[label=(\alph*)]
    \item If $v \cdot u = v \cdot w$ and $v \neq 0$, then $u = w$.
    \item If $u$ is a unit vector that is parallel to a nonzero vector $v$, then $u \cdot v = \pm \lVert v \rVert$.
\end{enumerate}

\medskip
\solution

\vspace{10em}

\hrulefill

\section*{Problem 5 (11.4 Cross product)}
Use determinant to find the cross product
\[
    i \times (i + j + k).
\]

\medskip
\solution

\vspace{6em}

\hrulefill

\section*{Problem 6 (11.4 Cross product)}
Let $u = \langle 1, 2, -3 \rangle$, $v = \langle -4, 1, 3 \rangle$. Find $u \times v$ and check that it is orthogonal to both $u$ and $v$.

\medskip
\solution

\vspace{10em}

\hrulefill

\section*{Problem 7 (11.4 Cross product)}
Find two units vectors that are orthogonal to both
\[
    u = -4 i + 3 j + k, \qquad v = 2 i + 4 k.
\]

\medskip
\solution

\vspace{10em}

\hrulefill

\section*{Problem 8 (11.4 Cross product)}
Determine whether the following points are co-planar: $P = (4, -3, 1)$, $Q = (6, -4, 7)$, $R = (1, 2, 2)$, and $S = (0, 1, 11)$.

\medskip
\solution

\vspace{10em}

\hrulefill

\end{document}
